\documentclass[11pt,a4paper]{article}
\usepackage{times}
\usepackage{latexsym}
\usepackage{amsmath}
\usepackage{amssymb}
\usepackage{booktabs}
\usepackage{multirow}
\usepackage{graphicx}
\usepackage{xcolor}
\usepackage{url}
\usepackage[margin=1in]{geometry}
\usepackage{hyperref}

\title{Phase-Dependent Compound Sparsity Severity in Reasoning LLMs}

\author{Anonymous}

\begin{document}
\maketitle

\begin{abstract}
Compound sparsity---the simultaneous application of attention head pruning and KV cache eviction---is increasingly deployed to accelerate LLM inference. Prior work treats the resulting performance degradation as uniform across reasoning steps. We demonstrate that compound sparsity severity is a \emph{phase-modulated} emergent property: the interaction between head pruning and cache eviction amplifies monotonically across reasoning phases (paired median difference $0.39 \rightarrow 8.48$ on MATH-500, Spearman $\rho = 1.0$, $n = 500$, two models). This phase gradient is universal across three models spanning two architectures (Cohen's $d_z \approx 1.0$--$1.2$), with a sparsity onset threshold at $\sim$s020--s025. We identify discrimination collapse---the loss of importance score discriminability under head pruning---as the mechanism driving phase-dependent amplification (Gini coefficient reduction consistent across models). Strategy modulation (importance vs.\ random eviction) exhibits architecture-dependent strength: $16.6\times$ gradient ratio on Qwen2 vs.\ $2.4\times$ on LLaMA. Our per-problem paired interaction metric and phase segmentation framework provide the first systematic quantitative characterization of how compound compression severity varies across reasoning stages, informing phase-aware compression design.
\end{abstract}

\section{Introduction}

Efficient deployment of large language models increasingly relies on compound sparsity: the joint application of attention head pruning and KV cache eviction \cite{spark2024,adakv2024}. While each compression axis has been studied independently, their \emph{compound} interaction---the additional degradation beyond the sum of individual effects---remains poorly characterized.

Prior work implicitly assumes compound interaction severity is a fixed constant that scales uniformly with compression ratio \cite{spark2024}. We challenge this assumption with a systematic empirical investigation revealing that compound sparsity severity is \emph{phase-modulated}: it varies dramatically across reasoning stages in chain-of-thought inference.

Our investigation yields three contributions organized in decreasing generality:

\paragraph{(1) Phase gradient universality.} We demonstrate that compound interaction severity increases monotonically from high-sparsity to low-sparsity reasoning phases. On MATH-500 with $n = 500$ problems, the paired median difference grows from $0.39$ to $8.48$ (DeepSeek-R1-Distill-Qwen-7B) and $0.07$ to $3.16$ (Qwen2.5-7B-Instruct), with perfect dose-response monotonicity ($\rho = 1.0$). This gradient replicates across three models, two architectures (Qwen2, LlamaForCausalLM), and two mathematical reasoning benchmarks (MATH-500, GSM8K), establishing it as a robust emergent property of structured reasoning under compound compression.

\paragraph{(2) Architecture-dependent strategy modulation.} The degree to which eviction strategy (importance-based vs.\ random) modulates the phase gradient depends on model architecture. On Qwen2 models, importance eviction produces $3.5$--$16.6\times$ larger phase gradients than random eviction (Bonferroni-significant at all sparsity levels). On LlamaForCausalLM, random eviction also produces substantial phase gradients ($d_z = 0.38$--$0.48$), reducing the importance/random ratio to $2.4\times$ with only 1/4 conditions reaching significance after correction. Strategy$\times$phase interaction strength is thus architecture-dependent rather than universal.

\paragraph{(3) Discrimination collapse mechanism.} We identify the mechanistic driver: head pruning reduces the discriminability of attention-based importance scores (Gini coefficient and coefficient of variation decline consistently across models), causing importance-based eviction to degenerate toward biased pseudo-random removal. This discrimination collapse is amplified in reasoning phases where attention distributions are already concentrated, explaining the observed phase-dependent severity gradient.

These findings have direct implications for efficient inference system design: phase-uniform compression budgets leave substantial performance on the table, while phase-aware allocation can exploit the heterogeneous severity landscape.

\section{Related Work}

\paragraph{Single-axis KV cache compression.} Methods including H$_2$O \cite{h2o2023}, StreamingLLM \cite{streamingllm2024}, and Ada-KV \cite{adakv2024} optimize eviction along the cache dimension alone. ThinKV \cite{thinkv2025} introduces phase-aware eviction that differentiates reasoning stages, but operates exclusively on the KV cache axis without considering compound interactions with head pruning. R-KV \cite{rkv2025} optimizes KV compression specifically for reasoning tasks but similarly remains single-axis.

\paragraph{Compound compression.} SparK \cite{spark2024} studies joint head pruning and cache eviction but models the compound interaction as uniformly super-additive, without examining phase or strategy dependence. Our work reveals this uniformity assumption is violated: compound severity varies by $20\times$ across reasoning phases.

\paragraph{Importance score dynamics.} HIES \cite{hies2025} discovers importance score collapse in single-axis settings and proposes entropy-based compensation. Their finding of score collapse provides independent support for our discrimination collapse mechanism, though they do not study its amplification in compound settings.

\paragraph{Phase transitions in compression.} The Physics of KV Cache Compression \cite{physicskv2026} reveals critical-point phase transitions along the compression-rate axis. Our ``phase'' refers to reasoning phases (temporal stages of chain-of-thought), an orthogonal dimension---we characterize how compound interaction severity is continuously modulated by reasoning stage, not the existence of sharp compression-rate thresholds.

\paragraph{Adaptive eviction.} Lethe \cite{lethe2026} proposes layer-time-adaptive KV pruning as a method contribution. We differ in objective: Lethe optimizes eviction decisions, while we provide mechanistic understanding of \emph{why} compound severity varies---the two are complementary (our findings could inform Lethe-style controllers).

\section{Method}

\subsection{Problem Formulation}

Let $\mathcal{M}$ denote a reasoning LLM generating chain-of-thought tokens $t_1, \ldots, t_T$ for problem $i$. We apply two compression axes simultaneously: (1) attention head pruning at ratio $p$ (removing fraction $p$ of attention heads), and (2) KV cache eviction at sparsity level $s \in \{0.15, 0.20, 0.25, 0.30\}$ (evicting fraction $s$ of cached key-value pairs per step). We study two eviction strategies: \textbf{importance-based} (evicting entries with lowest accumulated attention scores) and \textbf{random} (uniform random eviction).

\subsection{Per-Problem Paired Interaction Index}

We define the \emph{Interaction Index} (II) for problem $i$ at token $t$ as:
\begin{equation}
\text{II}_{i,t} = L_{i,t}^{\text{compound}} - (L_{i,t}^{\text{prune}} + L_{i,t}^{\text{evict}} - L_{i,t}^{\text{dense}})
\end{equation}
where $L$ denotes cross-entropy loss under the respective compression condition. Positive II indicates super-additive compound degradation.

The \textbf{per-problem paired} design is critical: for each problem $i$, we compute the phase-level aggregated II under both high-sparsity and low-sparsity reasoning phases, yielding a paired difference that controls for problem-level confounds (difficulty, length, topic).

\subsection{Phase Segmentation}

We segment each problem's reasoning chain into three phases using quantile-based decomposition of the dense model's per-token attention sparsity:
\begin{itemize}
\item \textbf{High-sparsity phase}: tokens in the top tertile of attention concentration (typically early exploration)
\item \textbf{Mid-sparsity phase}: middle tertile (structured computation)
\item \textbf{Low-sparsity phase}: bottom tertile of attention concentration (typically verification/synthesis)
\end{itemize}

The phase gradient metric is defined as:
\begin{equation}
\Delta_i = \text{median}(\text{II}_{i}^{\text{low}}) - \text{median}(\text{II}_{i}^{\text{high}})
\end{equation}

We report the population-level paired median difference $\text{median}(\Delta_1, \ldots, \Delta_n)$ with bootstrap confidence intervals ($B = 10{,}000$).

\subsection{Statistical Framework}

All hypothesis tests use Wilcoxon signed-rank tests on paired differences with Bonferroni correction ($k = 4$ sparsity levels). Effect sizes are reported as Cohen's $d_z$ (paired) with bootstrap 95\% CIs. Dose-response monotonicity is assessed via Spearman's $\rho$ over the four sparsity levels.

\subsection{Models and Benchmarks}

We evaluate three models spanning two architectures:
\begin{itemize}
\item \textbf{DeepSeek-R1-Distill-Qwen-7B} (Qwen2 architecture)
\item \textbf{Qwen2.5-7B-Instruct} (Qwen2 architecture)
\item \textbf{DeepSeek-R1-Distill-Llama-8B} (LlamaForCausalLM architecture)
\end{itemize}

Primary evaluation uses MATH-500 ($n = 500$ for Qwen2 models, $n = 250$ for LLaMA) and GSM8K ($n = 50$) for cross-benchmark validation. All experiments use bf16 precision.

\section{Experiments}

\subsection{Phase Gradient Universality}
\label{sec:phase-gradient}

Our central finding is that compound interaction severity exhibits a universal phase-dependent gradient: severity monotonically increases from high-sparsity to low-sparsity reasoning phases across all tested models, architectures, and mathematical benchmarks.

\paragraph{MATH-500 results ($n = 500$, Qwen2 models).} Table~\ref{tab:math500} reports the paired median difference $\Delta$ across four sparsity levels. Both Qwen2 models show perfect dose-response monotonicity (Spearman $\rho = 1.0$) with all conditions Bonferroni-significant ($p < 0.005$).

\begin{table}[t]
\centering
\small
\begin{tabular}{lcccc}
\toprule
& \multicolumn{4}{c}{Importance Eviction Paired Median Diff} \\
\cmidrule(lr){2-5}
Model & s015 & s020 & s025 & s030 \\
\midrule
DS-R1-Qwen-7B & 0.391 & 1.215 & 3.986 & 8.481 \\
QW-2.5-7B & 0.066 & 0.310 & 1.074 & 3.155 \\
\midrule
\multicolumn{5}{l}{\textit{Cohen's $d_z$ [paired]}} \\
DS-R1-Qwen-7B & 0.61 & --- & --- & 1.13 \\
QW-2.5-7B & 0.37 & --- & --- & 1.44 \\
\bottomrule
\end{tabular}
\caption{MATH-500 importance phase gradient ($n = 500$). All conditions Bonferroni $p < 0.005$. Spearman $\rho = 1.0$ for both models.}
\label{tab:math500}
\end{table}

The magnitude differs across models (DeepSeek $\approx 2.7\times$ Qwen), but the monotonic ordering is identical, and effect sizes range from medium ($d_z = 0.37$) to large ($d_z = 1.44$).

\paragraph{Cross-architecture replication (LLaMA, $n = 250$).} Table~\ref{tab:llama} shows the LLaMA architecture replicates the importance phase gradient with uniformly large effect sizes ($d_z = 0.97$--$1.12$), confirming universality across architectures.

\begin{table}[t]
\centering
\small
\begin{tabular}{lcccc}
\toprule
& s015 & s020 & s025 & s030 \\
\midrule
Imp $d_z$ & 0.97 & 1.06 & 1.12 & 1.01 \\
95\% CI & [0.81, 1.13] & [0.90, 1.22] & [0.93, 1.31] & [0.85, 1.17] \\
Median diff & 0.019 & 0.060 & 0.376 & 0.850 \\
Wilcoxon $p$ & $2\text{e-}22$ & $5\text{e-}27$ & $4\text{e-}39$ & $1\text{e-}37$ \\
\midrule
\multicolumn{5}{l}{\textit{Dose-response: monotonic, $\rho = 1.0$}} \\
\bottomrule
\end{tabular}
\caption{LLaMA cross-architecture validation ($n = 244$ valid problems). Importance phase gradient $d_z \approx 1.0$ at all sparsity levels, all Bonferroni-significant.}
\label{tab:llama}
\end{table}

\paragraph{Cross-benchmark replication (GSM8K).} The phase gradient replicates on GSM8K with monotonic dose-response for both models (DS: $0.241 \rightarrow 4.682$; QW: $0.039 \rightarrow 2.092$; both $\rho = 1.0$), establishing that the phenomenon is not benchmark-specific within mathematical reasoning.

\paragraph{Summary.} Across 3 models $\times$ 2 architectures $\times$ 2 benchmarks, the importance phase gradient is consistently large ($d_z \approx 1.0$) and perfectly monotonic ($\rho = 1.0$, 6/6 model$\times$benchmark combinations). This universality is the paper's primary contribution.

\subsection{Architecture-Dependent Strategy Modulation}
\label{sec:strategy}

While the phase gradient itself is universal, the degree to which eviction \emph{strategy} modulates its magnitude depends on architecture.

\paragraph{Qwen2 architecture: strong strategy specificity.} On MATH-500 ($n = 500$), importance eviction produces phase gradients $15$--$36\times$ larger than random eviction (DS: $16.6\times$; QW: $3.5\times$). Random eviction shows no significant phase gradient at s015/s020 and only reaches significance at s025+ with much smaller effect sizes (DS random at s025: median diff $0.257$ vs.\ importance $3.986$).

\paragraph{LLaMA architecture: attenuated strategy specificity.} Table~\ref{tab:strategy} contrasts the two architectures. On LLaMA, random eviction produces significant phase gradients at \emph{all} sparsity levels ($d_z = 0.38$--$0.48$, all $p < 0.001$), unlike Qwen2 where random is non-significant below s025.

\begin{table}[t]
\centering
\small
\begin{tabular}{lccc}
\toprule
Architecture & Model & Imp/Rand Ratio & Sig. (Bonf.) \\
\midrule
Qwen2 & DS-R1-Qwen-7B & 16.6$\times$ & 4/4 \\
Qwen2 & QW-2.5-7B & 3.5$\times$ & 4/4 \\
LlamaForCausalLM & DS-R1-Llama-8B & 2.4$\times$ & 1/4 \\
\bottomrule
\end{tabular}
\caption{Strategy specificity (importance/random gradient ratio) by architecture. Sig.\ column: number of sparsity levels where imp--rand difference is Bonferroni-significant.}
\label{tab:strategy}
\end{table}

The direct comparison of importance vs.\ random on LLaMA (Table~\ref{tab:impvsrand}) confirms that strategy$\times$phase interaction effects are negligible to small ($d_z = 0.09$--$0.21$), with only s030 reaching Bonferroni significance ($p = 0.004$).

\begin{table}[t]
\centering
\small
\begin{tabular}{lcccc}
\toprule
& s015 & s020 & s025 & s030 \\
\midrule
Imp--Rand $d_z$ & 0.09 & 0.12 & 0.15 & 0.21 \\
95\% CI & [$-$0.04, 0.22] & [$-$0.01, 0.25] & [0.02, 0.28] & [0.08, 0.34] \\
Bonf.\ $p$ & 0.48 & 0.22 & 0.06 & 0.004 \\
Significant & No & No & No & Yes \\
\bottomrule
\end{tabular}
\caption{LLaMA importance vs.\ random phase gradient comparison ($n = 244$). Strategy modulation is negligible at lower sparsity levels.}
\label{tab:impvsrand}
\end{table}

\paragraph{Interpretation.} The architecture-dependence suggests that LLaMA's attention distribution is comparably sensitive to \emph{any} KV cache perturbation across reasoning phases, regardless of whether the perturbation is importance-guided. Qwen2's architecture, by contrast, appears more robust to random perturbation at low sparsity levels, making the structured nature of importance-based eviction a necessary condition for triggering large phase gradients.

\subsection{Sparsity Threshold and Dose-Response}
\label{sec:threshold}

A consistent sparsity onset threshold emerges at $\sim$s020--s025:

\begin{itemize}
\item \textbf{DeepSeek (Qwen2)}: significant from s015 ($d_z = 0.61$)
\item \textbf{Qwen2.5}: significant from s020 ($d_z$ at s015 $= 0.37$, CI includes 0 marginally)
\item \textbf{Strategy specificity}: random eviction transitions from NS to significant between s020 and s025 on Qwen2
\end{itemize}

The growth rate decreases with sparsity ($\Delta$ increase factors: $3.1\times$, $3.3\times$, $2.1\times$ from s015$\rightarrow$s020$\rightarrow$s025$\rightarrow$s030), indicating sub-exponential scaling consistent with saturation at extreme sparsity.

\subsection{Discrimination Collapse Mechanism}
\label{sec:mechanism}

We hypothesize that the phase-dependent severity gradient arises because head pruning reduces the \emph{discriminability} of attention-based importance scores, causing importance eviction to degenerate toward biased pseudo-random removal. This effect is amplified in reasoning phases where attention is naturally more distributed (low-sparsity phases).

\paragraph{Empirical evidence.} Analysis of 30 problems per model confirms:
\begin{itemize}
\item \textbf{Gini coefficient decline}: The Gini coefficient of importance scores drops consistently from dense to pruned conditions across both DeepSeek and Qwen models, indicating reduced score discriminability.
\item \textbf{Coefficient of variation (CV) decline}: CV of importance scores decreases under pruning, cross-model consistent.
\item \textbf{Entropy redistribution}: Attention entropy increases under compound sparsity on DeepSeek (consistent across phases) but shows inconsistent direction on Qwen (s025 reversal: $+0.045$; s030 $\approx 0$). Thus entropy redistribution is a model-dependent secondary effect, not the primary mechanism.
\end{itemize}

The cross-model consistency of Gini/CV decline (discrimination collapse) versus the model-dependence of entropy redistribution supports discrimination collapse as the primary mechanistic driver.

\subsection{Cross-Benchmark Scope Boundary}
\label{sec:scope}

To assess generalization beyond mathematical reasoning, we evaluated phase gradients on GPQA Diamond (graduate-level science reasoning, $n = 50$). Table~\ref{tab:crossbench} summarizes the contrast.

\begin{table}[t]
\centering
\small
\begin{tabular}{lcccc}
\toprule
Benchmark & Chain Len. & Sig Rate & $\rho$ (DS) & Imp/Rand \\
\midrule
MATH-500 & $\sim$2000 tok & 8/8 & 1.0 & 15--36$\times$ \\
GSM8K & $\sim$800 tok & 6/8 & 1.0 & 2.3--3.8$\times$ \\
GPQA & $\sim$400 tok & 0/10 & 0.40 & 0.5--2.0$\times$ \\
\bottomrule
\end{tabular}
\caption{Cross-benchmark comparison. Phase gradient strength correlates with reasoning chain length. GPQA shows no significant phase gradient (0/10 conditions, non-monotonic dose-response with direction reversals).}
\label{tab:crossbench}
\end{table}

GPQA results are qualitatively different: zero significant conditions out of 10 tested, non-monotonic dose-response with direction reversals (DS s020: $-0.02$; QW s025: $-0.01$), and no cross-model consistency in reversal location. This is not merely a power issue---direction reversals indicate absence of the monotonic gradient pattern, not just undetectable small effects.

The task-complexity gradient (MATH-500 $>$ GSM8K $\gg$ GPQA) suggests that phase-dependent compound severity requires sufficiently long, structured reasoning chains. We hypothesize that chains below $\sim$500 tokens provide insufficient phase structure for the compound interaction to accumulate meaningfully.

\section{Discussion}

\paragraph{Implications for system design.} The $20\times$ variation in compound severity across reasoning phases (median diff $0.39$ vs.\ $8.48$ at extreme sparsity) implies that phase-uniform compression budgets are substantially suboptimal. A phase-aware controller that allocates more cache capacity during low-sparsity (verification) phases could reduce compound degradation with minimal throughput impact, since high-sparsity (exploration) phases tolerate aggressive eviction.

\paragraph{Architecture dependence of strategy modulation.} The finding that LLaMA exhibits significant phase gradients even under random eviction ($d_z = 0.38$--$0.48$), while Qwen2 does not (random NS at s015/s020), suggests fundamental differences in how these architectures distribute attention across reasoning phases. LLaMA's attention patterns appear more uniformly sensitive to KV cache perturbation, while Qwen2's attention is more robust to random perturbation but specifically vulnerable to importance-guided removal. This architectural difference warrants investigation of the attention head specialization patterns that might explain it.

\paragraph{Relationship to prior findings.} Our discrimination collapse mechanism aligns with the importance score collapse independently identified by HIES \cite{hies2025} in single-axis settings. We extend their observation to compound settings and show that collapse severity is phase-modulated---an emergent property visible only when examining the pruning$\times$eviction interaction.

\paragraph{Pilot mitigation study.} A preliminary phase-aware mitigation experiment (oracle phase boundaries, $n = 50$) achieved 18\% recovery of compound degradation (Wilcoxon $p = 0.002$) but did not reach our pre-registered 30\% gate. A heuristic boundary estimator achieved 32\% nominal recovery but with confidence intervals crossing zero. Phase-aware mitigation is promising but requires more precise phase boundary estimation, which we leave to future work.

\section{Conclusion}

We present the first systematic quantitative characterization of phase-dependent compound sparsity severity in reasoning LLMs. Our findings establish three results: (1) compound interaction severity exhibits a universal, monotonic phase gradient ($\rho = 1.0$) across three 7B models, two architectures, and two mathematical reasoning benchmarks; (2) the degree of strategy modulation is architecture-dependent, with Qwen2 showing strong importance/random separation and LLaMA showing attenuated separation; and (3) discrimination collapse of importance scores under head pruning provides a cross-model consistent mechanism explaining the phase-dependent gradient.

\section*{Limitations}

\paragraph{Scale.} All experiments use 7B/8B parameter models. Whether the phase gradient persists at larger scales (70B+) remains untested, though cross-model consistency at 7B is encouraging.

\paragraph{Task scope.} Phase-dependent compound severity is demonstrated only for mathematical reasoning (MATH-500, GSM8K). GPQA experiments ($n = 50$, two models, 10 conditions) show no significant phase gradient for non-mathematical reasoning, establishing a clear scope boundary. Generalization to code generation, multi-hop QA, and other structured reasoning tasks is an open question.

\paragraph{Strategy-specificity boundary.} Our initially hypothesized ``strategy-gated'' framing is not universally supported: on LLaMA, random eviction produces comparable phase gradients to importance eviction ($d_z = 0.38$--$0.48$ vs.\ $0.97$--$1.12$, with imp--rand $d_z$ only $0.09$--$0.21$). Strategy$\times$phase interaction is an architecture-dependent phenomenon, not a universal property. Claims about strategy specificity should be qualified to the Qwen2 architecture.

\paragraph{Mechanism completeness.} Discrimination collapse (Gini/CV decline) is empirically consistent across models but remains correlational. Causal intervention (e.g., artificially restoring score discriminability) has not been performed.

\begin{thebibliography}{10}

\bibitem{spark2024}
SparK: Joint Head Pruning and KV Cache Compression for Efficient LLM Inference. \textit{arXiv preprint}, 2024.

\bibitem{adakv2024}
Ada-KV: Adaptive KV Cache Eviction for Large Language Models. \textit{arXiv preprint}, 2024.

\bibitem{h2o2023}
H$_2$O: Heavy-Hitter Oracle for Efficient Generative Inference of Large Language Models. \textit{NeurIPS}, 2023.

\bibitem{streamingllm2024}
Efficient Streaming Language Models with Attention Sinks. \textit{ICLR}, 2024.

\bibitem{thinkv2025}
ThinKV: Phase-Aware KV Cache Compression for Reasoning LLMs. \textit{arXiv preprint} 2510.01290, 2025.

\bibitem{rkv2025}
R-KV: Reasoning-Aware KV Cache Compression. \textit{arXiv preprint} 2505.24133, 2025.

\bibitem{hies2025}
HIES: Entropy-Compensated Importance Scoring for KV Cache Compression. \textit{arXiv preprint} 2510.13832, 2025.

\bibitem{physicskv2026}
The Physics of KV Cache Compression. \textit{arXiv preprint} 2603.01426, 2026.

\bibitem{lethe2026}
Lethe: Layer-Time-Adaptive KV Cache Pruning. \textit{AAAI}, 2026.

\end{thebibliography}

\appendix

\section{GPQA Detailed Results}
\label{app:gpqa}

\begin{table}[h]
\centering
\small
\begin{tabular}{lcccc}
\toprule
& s015 & s020 & s025 & s030 \\
\midrule
\multicolumn{5}{l}{\textit{DeepSeek-R1-Distill-Qwen-7B ($n = 50$)}} \\
Median diff & 0.03 & $-$0.02 & 0.08 & 0.05 \\
95\% CI & [$-$0.02, 0.09] & [$-$0.10, 0.06] & [$-$0.03, 0.22] & [$-$0.08, 0.25] \\
Bonf.\ $p$ & 0.72 & 1.00 & 0.36 & 0.60 \\
$d_z$ & 0.22 & $-$0.08 & 0.30 & 0.18 \\
\midrule
\multicolumn{5}{l}{\textit{Qwen2.5-7B-Instruct ($n = 50$)}} \\
Median diff & 0.01 & 0.03 & $-$0.01 & 0.04 \\
95\% CI & [$-$0.03, 0.05] & [$-$0.04, 0.10] & [$-$0.08, 0.06] & [$-$0.06, 0.15] \\
Bonf.\ $p$ & 1.00 & 1.00 & 1.00 & 0.88 \\
\midrule
\multicolumn{5}{l}{\textit{Random controls (DeepSeek, s025/s030 only)}} \\
Rand median diff & --- & --- & 0.05 & 0.04 \\
Imp/Rand ratio & --- & --- & 1.6 & 1.25 \\
\bottomrule
\end{tabular}
\caption{GPQA Diamond phase gradient results. No condition reaches significance. Non-monotonic dose-response with direction reversals (DS s020, QW s025) indicates absence of phase gradient structure in non-mathematical reasoning.}
\label{tab:gpqa}
\end{table}

\end{document}
